% Options for packages loaded elsewhere
\PassOptionsToPackage{unicode}{hyperref}
\PassOptionsToPackage{hyphens}{url}
\PassOptionsToPackage{dvipsnames,svgnames,x11names}{xcolor}
\documentclass[
  russian,
  11pt,
  a4paper,
]{article}
\usepackage{xcolor}
\usepackage[top=22mm,bottom=22mm,left=22mm,right=22mm]{geometry}
\usepackage{amsmath,amssymb}
\setcounter{secnumdepth}{-\maxdimen} % remove section numbering
\usepackage{iftex}
\ifPDFTeX
  \usepackage[T1]{fontenc}
  \usepackage[utf8]{inputenc}
  \usepackage{textcomp} % provide euro and other symbols
\else % if luatex or xetex
  \usepackage{unicode-math} % this also loads fontspec
  \defaultfontfeatures{Scale=MatchLowercase}
  \defaultfontfeatures[\rmfamily]{Ligatures=TeX,Scale=1}
\fi
\usepackage{lmodern}
\ifPDFTeX\else
  % xetex/luatex font selection
  \setmainfont[]{Georgia}
  \setsansfont[]{Helvetica Neue}
  \setmonofont[]{Andale Mono}
\fi
% Use upquote if available, for straight quotes in verbatim environments
\IfFileExists{upquote.sty}{\usepackage{upquote}}{}
\IfFileExists{microtype.sty}{% use microtype if available
  \usepackage[]{microtype}
  \UseMicrotypeSet[protrusion]{basicmath} % disable protrusion for tt fonts
}{}
\makeatletter
\@ifundefined{KOMAClassName}{% if non-KOMA class
  \IfFileExists{parskip.sty}{%
    \usepackage{parskip}
  }{% else
    \setlength{\parindent}{0pt}
    \setlength{\parskip}{6pt plus 2pt minus 1pt}}
}{% if KOMA class
  \KOMAoptions{parskip=half}}
\makeatother
\ifLuaTeX
\usepackage[bidi=basic,shorthands=off,provide=*]{babel}
\else
\usepackage[bidi=default,shorthands=off,provide=*]{babel}
\fi
\ifPDFTeX
\else
\babelfont{rm}[]{Georgia}
\fi
\ifLuaTeX
  \usepackage{selnolig} % disable illegal ligatures
\fi
\setlength{\emergencystretch}{3em} % prevent overfull lines
\providecommand{\tightlist}{%
  \setlength{\itemsep}{0pt}\setlength{\parskip}{0pt}}
\usepackage{bookmark}
\IfFileExists{xurl.sty}{\usepackage{xurl}}{} % add URL line breaks if available
\urlstyle{same}
% fallback for those not using the hyperref driver hyperxmp:
\makeatletter
\@ifundefined{xmpquote}{\newcommand{\xmpquote}[1]{#1}}{}
\makeatother
\hypersetup{
  pdftitle={Планетарное соглашение о мире{,} безопасности{,} правах и общем знании},
  pdflang={ru-RU},
  colorlinks=true,
  linkcolor={blue},
  filecolor={Maroon},
  citecolor={Blue},
  urlcolor={blue},
  pdfcreator={LaTeX via pandoc}}

\title{Планетарное соглашение о мире, безопасности, правах и общем
знании}
\usepackage{etoolbox}
\makeatletter
\providecommand{\subtitle}[1]{% add subtitle to \maketitle
  \apptocmd{\@title}{\par {\large #1 \par}}{}{}
}
\makeatother
\subtitle{Проект рамочного договора}
\author{}
\date{14 августа 2026 года}

\begin{document}
\maketitle

\thispagestyle{empty}\clearpage

\renewcommand*\contentsname{Содержание}
{
\setcounter{tocdepth}{3}
\tableofcontents
}
Этот проект рамочного договора соединяет предотвращение войны, защиту
людей и планеты, справедливое правосудие и доступ к общему знанию в
одной проверяемой архитектуре. Он исходит из того, что устойчивый мир
невозможен одновременно без ограничения силы, без устранения
технологической зависимости и без защиты достоинства, частной жизни и
безопасности каждого человека.

\begin{quote}
\textbf{Статус документа:} концептуальный проект для политического,
общественного, научного и правового обсуждения. Он не является
действующим международным договором или готовым текстом для подписания и
требует переговоров, юридико-лингвистической экспертизы, согласования с
применимым международным правом, институционального статута, проверочных
протоколов и принятия Сторонами.
\end{quote}

\subsection{Краткая
архитектура}\label{ux43aux440ux430ux442ux43aux430ux44f-ux430ux440ux445ux438ux442ux435ux43aux442ux443ux440ux430}

Соглашение учреждает международную организацию с ограниченной переданной
компетенцией, но не мировое правительство и не орган общей власти. Оно
задаёт общие обязательства, независимую проверку, судебные средства
защиты и распределённые институты, исчерпывающие полномочия которых
ограничены текстом договора и его ратифицированными протоколами.

Его четыре равнообязательных начала:

\begin{enumerate}
\def\labelenumi{\arabic{enumi}.}
\tightlist
\item
  отказ от агрессии, поэтапное разоружение и мирное разрешение споров;
\item
  защита гражданского населения, критической инфраструктуры, биосферы и
  космического пространства;
\item
  неотъемлемые права, ответственность за тяжкие преступления и
  восстановление пострадавших;
\item
  презумпция открытости общественно значимого знания, сохранение
  цивилизационной памяти и адресное управление доказуемым риском.
\end{enumerate}

Ключевая информационная формула договора различает факт, сохранение,
доступ и ответственность. Состоявшуюся широкую публикацию нельзя честно
отменить юридической фикцией, однако это не создаёт безусловного права
продолжать распространение, индексирование или усиление вреда.
Сохранение уникального знания или доказательства не означает всеобщего
доступа к нему, а ограничение доступа не даёт права уничтожить память.

\subsection{Преамбула}\label{ux43fux440ux435ux430ux43cux431ux443ux43bux430}

Стороны настоящего Соглашения,

признавая равное достоинство всех людей, право народов самостоятельно
определять своё будущее и ответственность человечества перед будущими
поколениями,

подтверждая цели и принципы Устава Организации Объединённых Наций,
включая суверенное равенство государств, неприменение силы и мирное
разрешение международных споров,

подтверждая обязательства, вытекающие из международного права прав
человека, международного гуманитарного права, международного уголовного
права и права беженцев,

осуждая агрессию, геноцид, преступления против человечности, военные
преступления, террор против гражданского населения, захват территории
силой и коллективное наказание,

признавая, что ядерное, химическое и биологическое оружие,
неконтролируемые автономные вооружения, атаки на критическую
инфраструктуру, климатические и биологические катастрофы, опасное
использование цифровых систем и разрушение общей памяти создают
взаимосвязанные планетарные риски,

считая, что наука, культура, образование, технологическое сотрудничество
и справедливый доступ к жизненно важным знаниям уменьшают зависимость,
неравенство и причины конфликтов,

признавая одновременно право каждого пользоваться достижениями науки,
право автора на признание и справедливое вознаграждение, свободу
научного поиска, право на частную жизнь и право на эффективное средство
защиты,

исходя из презумпции максимальной открытости общественно значимого
знания при узких, доказуемых, срочных и пересматриваемых ограничениях,
необходимых для защиты людей от серьёзного и разумно предсказуемого
вреда,

признавая ценность распределённых архивов, открытых стандартов,
проверяемого происхождения, множественных интерпретаций и отсутствия
единственного органа, способного произвольно определять, что
человечеству разрешено помнить,

подтверждая, что космическое пространство должно использоваться в
интересах мира, международного сотрудничества и всех народов независимо
от уровня их экономического и научного развития,

договорились о нижеследующем.

\subsection{Часть I. Общие
положения}\label{ux447ux430ux441ux442ux44c-i-ux43eux431ux449ux438ux435-ux43fux43eux43bux43eux436ux435ux43dux438ux44f}

\subsubsection{Статья 1. Цели
Соглашения}\label{ux441ux442ux430ux442ux44cux44f-1-ux446ux435ux43bux438-ux441ux43eux433ux43bux430ux448ux435ux43dux438ux44f}

\begin{enumerate}
\def\labelenumi{\arabic{enumi}.}
\tightlist
\item
  Целями настоящего Соглашения являются:

  \begin{enumerate}
  \def\labelenumii{\arabic{enumii}.}
  \tightlist
  \item
    предотвращение войны, агрессии и эскалации вооружённых конфликтов;
  \item
    поэтапное, взаимное и проверяемое сокращение вооружений и связанных
    с ними рисков;
  \item
    защита гражданского населения, природы, критической гражданской
    инфраструктуры и планетарных общих ресурсов;
  \item
    обеспечение правосудия, установление истины, возмещение вреда и
    гарантии неповторения;
  \item
    справедливое распространение благ науки и технологий;
  \item
    сохранение цивилизационной памяти и развитие распределённого фонда
    общего знания;
  \item
    создание мирных механизмов сотрудничества на Земле и за её
    пределами.
  \end{enumerate}
\item
  Соглашение является рамочным договором. Подробные режимы вооружений,
  проверки, знания, цифровой безопасности, космоса, институтов и
  переходного финансирования закрепляются протоколами, которые не
  ослабляют прямые обязательства Соглашения.
\item
  Все положения Соглашения обязательны для Сторон. Для особой защиты от
  ослабления по статье 43 ядро составляют статьи 4--27, 33, 37 и 39.
  Проверочный протокол обязателен для каждого принятого вооруженческого
  обязательства; отсутствие технического протокола не приостанавливает
  запрет агрессии, обязанность защищать гражданское население, сохранять
  доказательства и исполнять применимое международное право.
\item
  Проект протокола готовит межправительственная конференция с открытой
  экспертной процедурой, принимает тремя четвертями Сторон и передаёт на
  отдельную ратификацию. Он вступает в силу для принявших его Сторон
  через девяносто дней после достижения установленного в нём порога,
  который не может быть ниже двух третей Сторон. Обязанность принять
  связанный проверочный протокол входит в соответствующее вооруженческое
  обязательство; уклонение от ратификации подлежит процедуре соблюдения,
  но Ассамблея не может единолично создать новое материальное
  обязательство.
\end{enumerate}

\subsubsection{Статья 2. Сфера применения и
участники}\label{ux441ux442ux430ux442ux44cux44f-2-ux441ux444ux435ux440ux430-ux43fux440ux438ux43cux435ux43dux435ux43dux438ux44f-ux438-ux443ux447ux430ux441ux442ux43dux438ux43aux438}

\begin{enumerate}
\def\labelenumi{\arabic{enumi}.}
\tightlist
\item
  Соглашение обязательно для государств и региональных интеграционных
  организаций, ставших его Сторонами.
\item
  Стороны обеспечивают соблюдение применимых обязательств своими
  органами, вооружёнными силами, организациями, находящимися под их
  юрисдикцией или фактическим контролем, и лицами, действующими по их
  поручению.
\item
  Города, местные сообщества, коренные народы, научные учреждения,
  компании, профессиональные объединения, организации гражданского
  общества и будущие постоянные поселения за пределами Земли могут
  принимать декларации ассоциированного участия. Такое участие не
  заменяет ответственности Сторон.
\item
  Негосударственные вооружённые субъекты могут принять односторонние
  обязательства о соблюдении гуманитарных, мирных и проверочных
  положений. Принятие не придаёт им государственного статуса и не
  освобождает от ответственности.
\item
  Региональная интеграционная организация при присоединении перечисляет
  переданную ей компетенцию, распределение ответственности с
  государствами-членами и последующие изменения. Организация и её члены
  не осуществляют одно и то же право голоса одновременно.
\item
  Ассоциированное участие реализуется через применимое национальное
  право, специальное соглашение или отдельный протокол и само по себе не
  создаёт международной правосубъектности.
\end{enumerate}

\subsubsection{Статья 3.
Определения}\label{ux441ux442ux430ux442ux44cux44f-3-ux43eux43fux440ux435ux434ux435ux43bux435ux43dux438ux44f}

Для целей Соглашения:

\begin{enumerate}
\def\labelenumi{\arabic{enumi}.}
\tightlist
\item
  \textbf{планетарный мир} означает состояние, при котором споры
  разрешаются без агрессии и угрозы силой, а общие институты
  предотвращают эскалацию, защищают права и устраняют причины системного
  насилия;
\item
  \textbf{акт агрессии} означает применение вооружённой силы
  государством против суверенитета, территориальной целостности или
  политической независимости другого государства либо иным образом,
  несовместимым с Уставом Организации Объединённых Наций;
\item
  \textbf{преступление агрессии} означает планирование, подготовку,
  инициирование или осуществление лицом, способным фактически
  контролировать политические или военные действия государства, такого
  акта агрессии, который по характеру, тяжести и масштабу представляет
  явное нарушение Устава Организации Объединённых Наций;
\item
  \textbf{критическая гражданская инфраструктура} означает системы,
  прекращение работы которых создаёт существенную угрозу жизни, здоровью
  или базовым правам населения, включая водоснабжение, продовольствие,
  здравоохранение, энергетику, связь, транспорт, финансовую
  устойчивость, предупреждение катастроф и гражданские архивы;
\item
  \textbf{цивилизационная память} означает научные, культурные,
  исторические, образовательные, правовые и доказательственные
  материалы, сохранение которых необходимо для непрерывности знаний,
  подотчётности и развития человечества;
\item
  \textbf{публикация} означает предоставление информации неопределённому
  кругу лиц либо фактическое распространение среди такого числа
  независимых получателей, при котором контроль над всеми копиями
  объективно утрачен; намерение учитывается при ответственности, но не
  меняет факта распространения;
\item
  \textbf{факт состоявшейся публикации} означает проверяемое
  историческое событие распространения; его признание не означает
  одобрения материала и не создаёт иммунитета от ответственности;
\item
  \textbf{общедоступное знание} означает информацию, отнесённую статьёй
  20 или законным решением к открытому контуру и доступную на условиях,
  разрешающих указанные там способы использования;
\item
  \textbf{операционная информация высокого риска} означает достаточно
  конкретные, непосредственно применимые и актуальные сведения,
  распространение которых существенно облегчает причинение массового или
  иного тяжкого трансграничного вреда и риск которого нельзя разумно
  снизить менее ограничительным средством;
\item
  \textbf{защищаемые персональные или общинные сведения} означают
  сведения, доступ к которым затрагивает частную жизнь, безопасность,
  достоинство либо законно признанное коллективное право определяемых
  лиц и сообществ; решение учитывает легитимность представительства,
  внутреннее несогласие и общественную значимость раскрытия;
\item
  \textbf{общественно значимое раскрытие} означает сообщение сведений,
  которые лицо имеет разумные основания считать достоверно указывающими
  на тяжкое нарушение права, коррупцию, угрозу жизни, экологическую
  опасность, обман общества или злоупотребление властью;
\item
  \textbf{запечатанное хранение} означает сохранение целостной копии с
  ограниченным доступом, независимым хранителем, журналом обращений и
  определённой процедурой пересмотра;
\item
  \textbf{серьёзный и разумно предсказуемый вред} означает вред,
  вероятность и масштаб которого подтверждаются конкретными
  обстоятельствами, а не одной ссылкой на широкую категорию безопасности
  или двойного назначения.
\end{enumerate}

\subsubsection{Статья 4. Основные принципы и
толкование}\label{ux441ux442ux430ux442ux44cux44f-4-ux43eux441ux43dux43eux432ux43dux44bux435-ux43fux440ux438ux43dux446ux438ux43fux44b-ux438-ux442ux43eux43bux43aux43eux432ux430ux43dux438ux435}

\begin{enumerate}
\def\labelenumi{\arabic{enumi}.}
\tightlist
\item
  Стороны действуют добросовестно, с соблюдением требований законности,
  необходимости, соразмерности, недискриминации, прозрачности и
  подотчётности.
\item
  Ограничение права или доступа допускается только на основании
  опубликованного закона, ради легитимной цели и при отсутствии менее
  ограничительного средства. Индивидуальная мера и режим секретности
  имеют определённый срок; общая норма защиты периодически
  пересматривается.
\item
  Предосторожность не может служить бессрочной или бездоказательной
  секретности. Открытость не может служить оправданием конкретного,
  предсказуемого и предотвратимого тяжкого вреда.
\item
  Никакое положение не толкуется как разрешение агрессии, геноцида,
  преступления против человечности, военного преступления, пытки,
  рабства, насильственного исчезновения или дискриминации.
\item
  Соглашение не уменьшает защиту, предоставленную иными применимыми
  нормами. При нескольких допустимых толкованиях выбирается то, которое
  наиболее полно защищает мир и человеческое достоинство при наименьшем
  необоснованном ограничении знания.
\item
  Ничто в Соглашении не создаёт дополнительного основания для угрозы
  силой или её применения. Сила допускается исключительно в случаях и
  пределах, предусмотренных Уставом Организации Объединённых Наций;
  обязательства по Уставу имеют преимущество.
\item
  Простое добросовестное хранение, изучение, проверка, перевод или
  архивирование законного знания не создаёт ответственности без
  дополнительного противоправного действия, умысла либо небрежного
  причинения конкретного вреда.
\end{enumerate}

\subsection{Часть II. Мир и общая
безопасность}\label{ux447ux430ux441ux442ux44c-ii-ux43cux438ux440-ux438-ux43eux431ux449ux430ux44f-ux431ux435ux437ux43eux43fux430ux441ux43dux43eux441ux442ux44c}

\subsubsection{Статья 5. Неприменение
силы}\label{ux441ux442ux430ux442ux44cux44f-5-ux43dux435ux43fux440ux438ux43cux435ux43dux435ux43dux438ux435-ux441ux438ux43bux44b}

\begin{enumerate}
\def\labelenumi{\arabic{enumi}.}
\tightlist
\item
  Стороны воздерживаются от угрозы силой и её применения против
  территориальной целостности, политической независимости или населения
  любого государства и народа.
\item
  Территориальные приобретения, полученные в результате угрозы силой или
  её применения, не признаются законными.
\item
  Ничто в Соглашении не умаляет права на индивидуальную или коллективную
  самооборону при вооружённом нападении и в пределах международного
  права. Меры самообороны должны быть необходимыми, соразмерными,
  направленными на прекращение нападения и немедленно сообщаться Совету
  Безопасности Организации Объединённых Наций с дополнительным
  уведомлением Совета предотвращения конфликтов.
\item
  Стороны не организуют, не финансируют и не направляют вооружённые
  группы для действий, которые были бы запрещены самой Стороне.
\item
  Военная доктрина, закупки и развёртывание сил должны быть совместимы с
  целью снижения риска внезапного нападения и ошибочной эскалации.
\end{enumerate}

\subsubsection{Статья 6. Прекращение боевых действий и мирное разрешение
споров}\label{ux441ux442ux430ux442ux44cux44f-6-ux43fux440ux435ux43aux440ux430ux449ux435ux43dux438ux435-ux431ux43eux435ux432ux44bux445-ux434ux435ux439ux441ux442ux432ux438ux439-ux438-ux43cux438ux440ux43dux43eux435-ux440ux430ux437ux440ux435ux448ux435ux43dux438ux435-ux441ux43fux43eux440ux43eux432}

\begin{enumerate}
\def\labelenumi{\arabic{enumi}.}
\tightlist
\item
  Сторона вооружённого конфликта, для которой Соглашение вступило в
  силу, либо подписант, письменно заявивший о временном применении
  настоящей статьи, в течение двадцати четырёх часов прекращает
  наступательные операции и добросовестно согласовывает контролируемое
  прекращение огня, гуманитарные гарантии и разведение сил.
\item
  До разведения сил допускаются только необходимые и соразмерные
  оборонительные действия против продолжающегося или непосредственно
  предстоящего вооружённого нападения. Они не могут служить расширению
  территориального контроля или наказанию населения.
\item
  Инспекция либо согласованный нейтральный механизм немедленно наблюдает
  прекращение огня, устанавливает факты предполагаемых нарушений и
  публикует безопасный отчёт. Спор о нарушении не приостанавливает
  защиту гражданских и гуманитарные обязанности.
\item
  Стороны разрешают споры посредством консультаций, переговоров, добрых
  услуг, посредничества, установления фактов, арбитража или судебного
  разбирательства и применяют только мирные и законные невоенные
  временные меры.
\item
  Линии экстренной связи между политическими и военными центрами решений
  действуют непрерывно, регулярно проверяются и не отключаются как
  средство политического давления.
\item
  Посредничество не может навязывать потерпевшим отказ от неотъемлемых
  прав или безусловную амнистию за основные международные преступления.
\end{enumerate}

\subsubsection{Статья 7. Защита гражданского
населения}\label{ux441ux442ux430ux442ux44cux44f-7-ux437ux430ux449ux438ux442ux430-ux433ux440ux430ux436ux434ux430ux43dux441ux43aux43eux433ux43e-ux43dux430ux441ux435ux43bux435ux43dux438ux44f}

\begin{enumerate}
\def\labelenumi{\arabic{enumi}.}
\tightlist
\item
  Стороны всегда различают гражданских лиц и участников боевых действий,
  гражданские объекты и военные цели.
\item
  Запрещаются нападения на гражданское население, неизбирательные и
  несоразмерные атаки, голод как метод войны, насильственное перемещение
  без законных оснований, захват заложников и использование гражданских
  лиц как прикрытия.
\item
  Гражданский объект сохраняет защиту, пока он не вносит эффективный
  вклад в военные действия и его нейтрализация в данных обстоятельствах
  не даёт определённого военного преимущества. При сомнении объект
  считается гражданским. Выбирающий цель принимает все практически
  возможные меры проверки, предупреждения и уменьшения сопутствующего
  вреда.
\item
  Больницы, школы, объекты воды и продовольствия, установки, содержащие
  опасные силы, системы гражданской связи, гуманитарные службы,
  культурные ценности, архивы и средства проверки Соглашения пользуются
  усиленной защитой по применимому гуманитарному праву. Стороны не
  используют их для прикрытия военных целей; специальный протокол может
  установить функции, защита которых никогда не прекращается.
\item
  Стороны обеспечивают безопасный, быстрый и беспристрастный
  гуманитарный доступ, эвакуацию нуждающихся, обмен сведениями о
  пропавших и достойное обращение с умершими. Согласие на
  беспристрастную помощь не удерживается произвольно; разумный контроль
  маршрутов и безопасности не должен лишать помощь практического
  эффекта.
\item
  Все задержанные и военнопленные содержатся гуманно, регистрируются,
  получают связь с семьёй и доступ к независимому контролю. Стороны
  соблюдают невысылку при реальном риске преследования, пытки или иного
  тяжкого вреда, обеспечивают безопасное добровольное возвращение
  перемещённых лиц и принимают проверяемые программы разминирования,
  демобилизации и реинтеграции.
\item
  Цифровое отключение населения, если оно препятствует обращению за
  помощью, предупреждению опасности или документированию тяжких
  нарушений, не может применяться как общее средство контроля.
\end{enumerate}

\subsubsection{Статья 8. Оружие массового
уничтожения}\label{ux441ux442ux430ux442ux44cux44f-8-ux43eux440ux443ux436ux438ux435-ux43cux430ux441ux441ux43eux432ux43eux433ux43e-ux443ux43dux438ux447ux442ux43eux436ux435ux43dux438ux44f}

\begin{enumerate}
\def\labelenumi{\arabic{enumi}.}
\tightlist
\item
  Стороны не разрабатывают, не испытывают, не производят, не
  приобретают, не получают, не передают, не размещают, не хранят, не
  владеют, не применяют и не угрожают применением ядерного, химического
  или биологического оружия и не помогают в совершении таких действий.
\item
  Существующие запрещённые запасы, средства доставки, производственные
  возможности и программы декларируются и подлежат безопасной,
  необратимой и проверяемой ликвидации. До неё запас может временно
  находиться только под запечатанным контролем ради уничтожения; это
  узкое переходное исключение из запрета хранения и владения не
  распространяется на приобретение, передачу, размещение, угрозу или
  применение. В течение одного года принимается обязательный протокол с
  исходными уровнями, этапами, методами проверки и первым измеримым
  сокращением не позднее двух лет после его вступления в силу.
\item
  До ликвидации Стороны снижают боевую готовность и риск ошибочного
  пуска, обеспечивают разделённые полномочия, технические блокировки,
  безопасную связь и значимый человеческий контроль. Эти временные
  обязанности не означают признания законности владения, угрозы
  применением или применения.
\item
  Запрет не препятствует мирной химической, биологической, медицинской и
  энергетической деятельности, если её виды, количества, защитные меры и
  проверка совместимы с настоящей статьёй и применимым международным
  правом.
\item
  Инспекция раскрывает обществу выводы, методику и сведения, необходимые
  для подотчётности, но не публикует не нужные для этого операционные
  подробности. Такие подробности помещаются в защищённое приложение с
  управляемым доступом.
\end{enumerate}

\subsubsection{Статья 9. Обычные вооружения и меры
доверия}\label{ux441ux442ux430ux442ux44cux44f-9-ux43eux431ux44bux447ux43dux44bux435-ux432ux43eux43eux440ux443ux436ux435ux43dux438ux44f-ux438-ux43cux435ux440ux44b-ux434ux43eux432ux435ux440ux438ux44f}

\begin{enumerate}
\def\labelenumi{\arabic{enumi}.}
\tightlist
\item
  Стороны ограничивают вооружённые силы уровнем, достаточным для
  законной обороны и совместимых с Соглашением миротворческих задач. В
  течение одного года обязательный протокол устанавливает исходные
  уровни, численные и качественные пределы, пороги уведомления, этапы
  сокращения и методику проверки.
\item
  Крупные учения, перемещения сил и испытания систем, способные быть
  воспринятыми как подготовка нападения, заранее уведомляются через
  проверяемый механизм.
\item
  В приграничных районах создаются согласованные зоны сниженной военной
  активности, режимы наблюдения и процедуры предотвращения инцидентов.
\item
  Передача вооружений требует оценки риска их использования для
  агрессии, тяжких нарушений прав или обхода Соглашения. Высокий риск
  влечёт отказ в передаче.
\item
  Бюджеты, закупки и конечные владельцы вооружений раскрываются в
  объёме, достаточном для общественной и международной подотчётности.
\end{enumerate}

\subsubsection{Статья 10. Автономные вооружения и искусственный
интеллект}\label{ux441ux442ux430ux442ux44cux44f-10-ux430ux432ux442ux43eux43dux43eux43cux43dux44bux435-ux432ux43eux43eux440ux443ux436ux435ux43dux438ux44f-ux438-ux438ux441ux43aux443ux441ux441ux442ux432ux435ux43dux43dux44bux439-ux438ux43dux442ux435ux43bux43bux435ux43aux442}

\begin{enumerate}
\def\labelenumi{\arabic{enumi}.}
\tightlist
\item
  Решение о выборе человека целью, применении против него смертельной
  силы либо атаке на объект с разумно предсказуемыми смертельными
  последствиями не передаётся системе без значимого человеческого
  контроля.
\item
  Значимый человеческий контроль требует своевременной и достаточной
  информации, понимания цели, среды и вероятных последствий, реальной
  возможности отменить или прекратить действие и заранее определённого
  ответственного лица в командной цепочке. Система имеет проверяемые
  пределы применения, протокол испытаний, журнал существенных решений и
  возможность безопасного прекращения работы.
\item
  Запрещается развёртывание системы, поведение которой невозможно
  достаточно объяснить, ограничить или проверить применительно к
  предполагаемой среде.
\item
  Автоматизация не уменьшает ответственность государства, командира,
  оператора, разработчика или поставщика в пределах их фактической роли
  и применимого права.
\item
  Стороны обмениваются гражданскими методами проверки безопасности,
  данными об опасных отказах и мерами предотвращения неконтролируемой
  эскалации, не раскрывая персональные сведения и операционные ключи.
\end{enumerate}

\subsubsection{Статья 11. Кибермир и критическая
инфраструктура}\label{ux441ux442ux430ux442ux44cux44f-11-ux43aux438ux431ux435ux440ux43cux438ux440-ux438-ux43aux440ux438ux442ux438ux447ux435ux441ux43aux430ux44f-ux438ux43dux444ux440ux430ux441ux442ux440ux443ux43aux442ux443ux440ux430}

\begin{enumerate}
\def\labelenumi{\arabic{enumi}.}
\tightlist
\item
  Стороны не осуществляют и не поддерживают цифровые операции, основной
  или разумно предсказуемый результат которых состоит в причинении
  тяжкого вреда гражданскому населению.
\item
  Особо запрещаются атаки на исключительно гражданские функции
  здравоохранения, воды, продовольствия, ядерной безопасности, раннего
  предупреждения, гуманитарной связи, гражданских архивов и механизмов
  проверки Соглашения. Иная квалификация объекта определяется по статье
  7 и применимому международному гуманитарному праву.
\item
  Существенная уязвимость критической системы под юрисдикцией или
  контролем Стороны сообщается назначенному органу через защищённую
  процедуру координированного раскрытия. Исследователь получает
  безопасную гавань за добросовестное тестирование в разрешённых
  пределах; сведения об уже используемой атаке передаются без задержки.
  Ограничение публикации действует лишь в объёме и на срок, необходимые
  для исправления, и подтверждается независимым надзором.
\item
  Стороны не требуют универсальных скрытых доступов, системных
  ослаблений шифрования или тайных уязвимостей, которые создают
  несоразмерный общий риск.
\item
  Независимая техническая группа устанавливает факты цифровых инцидентов
  по проверяемой методике и отделяет степень уверенности от политической
  оценки.
\end{enumerate}

\subsubsection{Статья 12. Биологическая, экологическая и климатическая
безопасность}\label{ux441ux442ux430ux442ux44cux44f-12-ux431ux438ux43eux43bux43eux433ux438ux447ux435ux441ux43aux430ux44f-ux44dux43aux43eux43bux43eux433ux438ux447ux435ux441ux43aux430ux44f-ux438-ux43aux43bux438ux43cux430ux442ux438ux447ux435ux441ux43aux430ux44f-ux431ux435ux437ux43eux43fux430ux441ux43dux43eux441ux442ux44c}

\begin{enumerate}
\def\labelenumi{\arabic{enumi}.}
\tightlist
\item
  Стороны проявляют должную осмотрительность для предотвращения
  существенного трансграничного вреда здоровью, климату и экосистемам,
  проводят предварительную оценку воздействия, уведомляют потенциально
  затрагиваемые Стороны и устраняют разумно предсказуемый риск.
\item
  Данные о непосредственной угрозе жизни, эпидемии, загрязнении,
  радиационном выбросе и опасном изменении среды раскрываются без
  неоправданной задержки.
\item
  Исследования повышенного биологического риска проходят независимую
  оценку, многоуровневую физическую и информационную защиту,
  журналирование и международное уведомление в объёме, не создающем
  новой операционной угрозы.
\item
  Стороны сотрудничают в мониторинге, раннем предупреждении, медицинском
  ответе, восстановлении экосистем и справедливом распространении
  безопасных технологий.
\item
  Загрязнитель несёт обязанность прекратить вред, восстановить среду и
  компенсировать ущерб в соответствии с применимым правом и
  установленными фактами.
\end{enumerate}

\subsubsection{Статья 13. Космос и планетарные общие
ресурсы}\label{ux441ux442ux430ux442ux44cux44f-13-ux43aux43eux441ux43cux43eux441-ux438-ux43fux43bux430ux43dux435ux442ux430ux440ux43dux44bux435-ux43eux431ux449ux438ux435-ux440ux435ux441ux443ux440ux441ux44b}

\begin{enumerate}
\def\labelenumi{\arabic{enumi}.}
\tightlist
\item
  Исследование и использование космического пространства осуществляются
  мирно, в интересах всех стран и с учётом будущих поколений.
\item
  Запрещается размещать оружие массового уничтожения в космосе,
  создавать военные базы на небесных телах и применять силу для
  исключительного присвоения общих ресурсов.
\item
  Стороны предупреждают опасное загрязнение Земли и небесных тел,
  образование космического мусора и вредные помехи деятельности других
  участников.
\item
  Данные о планетарных угрозах, навигационной безопасности, космической
  погоде и опасных сближениях раскрываются своевременно и в открытых
  совместимых форматах.
\item
  Стороны разрешают деятельность своих государственных и
  негосударственных операторов и постоянно надзирают за ней. Добыча
  внеземных ресурсов не создаёт суверенитета над небесным телом; её
  экологические пределы, прозрачный учёт и режим выгод определяются
  отдельным протоколом без предрешения существующих разногласий до его
  принятия.
\item
  Экипажи и постоянные поселения оказывают взаимную помощь при угрозе
  жизни независимо от гражданства, политического статуса или
  принадлежности инфраструктуры.
\end{enumerate}

\subsubsection{Статья 14. Экономические меры и недопустимость
коллективного
наказания}\label{ux441ux442ux430ux442ux44cux44f-14-ux44dux43aux43eux43dux43eux43cux438ux447ux435ux441ux43aux438ux435-ux43cux435ux440ux44b-ux438-ux43dux435ux434ux43eux43fux443ux441ux442ux438ux43cux43eux441ux442ux44c-ux43aux43eux43bux43bux435ux43aux442ux438ux432ux43dux43eux433ux43e-ux43dux430ux43aux430ux437ux430ux43dux438ux44f}

\begin{enumerate}
\def\labelenumi{\arabic{enumi}.}
\tightlist
\item
  Экономические и финансовые меры реагирования должны иметь законную
  цель, быть адресными, соразмерными, временными, пересматриваемыми и
  допускающими отмену. Обязательства по решениям Совета Безопасности
  Организации Объединённых Наций имеют преимущество; Сторона всё равно
  использует доступные ей способы уменьшения гуманитарного вреда.
\item
  Они не могут намеренно лишать гражданское население продовольствия,
  воды, базовой медицины, средств жизнеобеспечения, связи с
  гуманитарными службами или безопасных образовательных знаний.
\item
  Каждая мера содержит гуманитарные исключения, процедуру исправления
  ошибочного включения и дату обязательного пересмотра.
\item
  Ограничения на гражданские технологии не применяются как средство
  закрепления постоянной зависимости. Реальный риск военного
  перенаправления устраняется адресными условиями, проверкой конечного
  использования и безопасным лицензированием.
\item
  Лица и организации, непосредственно ответственные за нарушение, не
  могут прикрываться запретом коллективного наказания, чтобы избежать
  индивидуальных мер и правосудия.
\end{enumerate}

\subsection{Часть III. Права, справедливость и
ответственность}\label{ux447ux430ux441ux442ux44c-iii-ux43fux440ux430ux432ux430-ux441ux43fux440ux430ux432ux435ux434ux43bux438ux432ux43eux441ux442ux44c-ux438-ux43eux442ux432ux435ux442ux441ux442ux432ux435ux43dux43dux43eux441ux442ux44c}

\subsubsection{Статья 15. Права и
недискриминация}\label{ux441ux442ux430ux442ux44cux44f-15-ux43fux440ux430ux432ux430-ux438-ux43dux435ux434ux438ux441ux43aux440ux438ux43cux438ux43dux430ux446ux438ux44f}

\begin{enumerate}
\def\labelenumi{\arabic{enumi}.}
\tightlist
\item
  Каждая Сторона уважает и обеспечивает права всех лиц под её
  юрисдикцией или фактическим контролем без дискриминации.
\item
  Меры отступления допускаются только при официально объявленном
  чрезвычайном положении, угрожающем жизни нации, строго в пределах
  необходимости, временно, без дискриминации и в соответствии с иными
  международными обязательствами. Сторона незамедлительно уведомляет
  депозитария о мерах, основаниях, территории и сроке. Отступление от
  прав, от которых применимое международное право его не допускает,
  запрещено; иные ограничения отвечают статье 4.
\item
  Никто не подвергается преследованию только за национальность,
  происхождение, язык, убеждения, религию, пол, гендер, инвалидность,
  возраст, имущественное положение или законное участие в мирной
  политической и научной деятельности.
\item
  Алгоритмические решения публичных органов, существенно затрагивающие
  права, должны быть объяснимыми, оспоримыми, проверяемыми на
  дискриминацию и подчинёнными эффективному человеческому пересмотру.
\item
  Лицо имеет право знать применимое к нему правило, существенные
  основания неблагоприятного решения и доступный порядок обжалования.
\end{enumerate}

\subsubsection{Статья 16. Потерпевшие, истина и
возмещение}\label{ux441ux442ux430ux442ux44cux44f-16-ux43fux43eux442ux435ux440ux43fux435ux432ux448ux438ux435-ux438ux441ux442ux438ux43dux430-ux438-ux432ux43eux437ux43cux435ux449ux435ux43dux438ux435}

\begin{enumerate}
\def\labelenumi{\arabic{enumi}.}
\tightlist
\item
  Потерпевшие имеют право на безопасность, участие, правду, доступ к
  правосудию, восстановление, компенсацию, реабилитацию и гарантии
  неповторения.
\item
  Расследования должны быть независимыми, беспристрастными,
  своевременными и способными установить индивидуальную и
  институциональную ответственность.
\item
  Государственные архивы, места захоронений, реестры пропавших, журналы
  приказов и иные доказательства тяжких нарушений сохраняются даже при
  смене власти или прекращении организации.
\item
  Публичное раскрытие учитывает безопасность потерпевших, детей,
  свидетелей и родственников. Защита личности не даёт права уничтожать
  запечатанную доказательственную копию.
\item
  Исправление ложной записи, связанное опровержение и признание вреда
  являются самостоятельными средствами защиты и не подменяют
  компенсацию.
\end{enumerate}

\subsubsection{Статья 17. Основные международные
преступления}\label{ux441ux442ux430ux442ux44cux44f-17-ux43eux441ux43dux43eux432ux43dux44bux435-ux43cux435ux436ux434ux443ux43dux430ux440ux43eux434ux43dux44bux435-ux43fux440ux435ux441ux442ux443ux43fux43bux435ux43dux438ux44f}

\begin{enumerate}
\def\labelenumi{\arabic{enumi}.}
\tightlist
\item
  Геноцид, преступления против человечности, военные преступления и
  преступление агрессии расследуются и преследуются независимыми
  национальными судами, Международным уголовным судом либо иным
  компетентным международным механизмом в пределах их юрисдикции и
  применимого права. Судебная палата настоящего Соглашения уголовным
  судом не является.
\item
  Уголовная ответственность основывается на законе, действовавшем во
  время деяния, личной вине и справедливом разбирательстве. Должность,
  приказ начальника, чрезвычайное положение или политическая
  целесообразность сами по себе не освобождают от уголовной
  ответственности, но вопросы иммунитета решаются по применимому
  международному и национальному праву.
\item
  Стороны устанавливают территориальную, персональную и иную
  согласованную юрисдикцию, сотрудничают в сохранении доказательств,
  правовой помощи, передаче дела, выдаче либо законном преследовании и
  защите свидетелей. Палата может передавать материалы компетентному
  органу, не предрешая виновность.
\item
  Никто не выдаётся и не передаётся при реальном риске пытки, смертной
  казни без надёжной гарантии её неприменения или явно несправедливого
  суда.
\item
  Амнистия и срок давности не могут применяться вопреки обязательству
  расследовать и преследовать основные международные преступления.
  Примирение основывается на установленных фактах, участии потерпевших и
  гарантиях неповторения, а не на обязательном забвении.
\end{enumerate}

\subsubsection{Статья 18. Народы, коренные сообщества и будущие
поколения}\label{ux441ux442ux430ux442ux44cux44f-18-ux43dux430ux440ux43eux434ux44b-ux43aux43eux440ux435ux43dux43dux44bux435-ux441ux43eux43eux431ux449ux435ux441ux442ux432ux430-ux438-ux431ux443ux434ux443ux449ux438ux435-ux43fux43eux43aux43eux43bux435ux43dux438ux44f}

\begin{enumerate}
\def\labelenumi{\arabic{enumi}.}
\tightlist
\item
  Ничто в Соглашении не толкуется как право насильственно изменять
  политический статус народа, уничтожать его язык, культуру или связь с
  землёй.
\item
  По решениям, затрагивающим права коренных народов, их территории,
  генетические ресурсы, сакральные материалы или традиционное знание,
  Стороны консультируются и сотрудничают с их легитимными
  представительными институтами с целью получить свободное,
  предварительное и осознанное согласие во всех случаях, предусмотренных
  применимым международным правом.
\item
  Открытость знания не отменяет атрибуцию, культурный контекст,
  справедливое распределение выгод и защиту сведений, раскрытие которых
  создаёт конкретный вред общине.
\item
  Крупные решения оцениваются с учётом мира, климата, биоразнообразия,
  долговой нагрузки, технологической зависимости и сохранности знания
  для будущих поколений.
\item
  Форум гражданского общества и будущих поколений вправе представлять
  заключения Ассамблее и Судебной палате.
\end{enumerate}

\subsection{Часть IV. Общее знание и цивилизационная
память}\label{ux447ux430ux441ux442ux44c-iv-ux43eux431ux449ux435ux435-ux437ux43dux430ux43dux438ux435-ux438-ux446ux438ux432ux438ux43bux438ux437ux430ux446ux438ux43eux43dux43dux430ux44f-ux43fux430ux43cux44fux442ux44c}

\subsubsection{Статья 19. Три контура обращения с
информацией}\label{ux441ux442ux430ux442ux44cux44f-19-ux442ux440ux438-ux43aux43eux43dux442ux443ux440ux430-ux43eux431ux440ux430ux449ux435ux43dux438ux44f-ux441-ux438ux43dux444ux43eux440ux43cux430ux446ux438ux435ux439}

\begin{enumerate}
\def\labelenumi{\arabic{enumi}.}
\tightlist
\item
  Соглашение устанавливает три контура:

  \begin{enumerate}
  \def\labelenumii{\arabic{enumii}.}
  \tightlist
  \item
    презумптивно открытое общественно значимое знание;
  \item
    контролируемый доступ к операционной информации высокого риска;
  \item
    защита персональных и общинных сведений.
  \end{enumerate}
\item
  Категория определяется по содержанию, контексту, актуальности и
  доказуемому риску, а не по происхождению, политической неудобности,
  новизне или общей ссылке на двойное назначение.
\item
  Один материал может содержать части разных контуров. В таком случае
  применяется отделение, редактирование или запечатанное приложение, а
  не автоматическое закрытие всего материала.
\item
  Лицо, требующее ограничения, несёт бремя показать правовое основание,
  конкретный риск, срок и отсутствие менее ограничительного средства.
\item
  Назначенный национальный орган может ввести узкое предварительное
  ограничение не более чем на семьдесят два часа. Комиссия независимо
  оценивает его и даёт рекомендацию; продление требует решения Судебной
  палаты или компетентного национального суда после проверки по
  приложению Б.
\item
  Решение подлежит независимому обжалованию. Секретное обоснование не
  может быть единственным основанием ограничения, а затронутой стороне
  сообщается безопасное резюме, достаточное для реального оспаривания.
\end{enumerate}

\subsubsection{Статья 20. Презумптивно открытое
знание}\label{ux441ux442ux430ux442ux44cux44f-20-ux43fux440ux435ux437ux443ux43cux43fux442ux438ux432ux43dux43e-ux43eux442ux43aux440ux44bux442ux43eux435-ux437ux43dux430ux43dux438ux435}

\begin{enumerate}
\def\labelenumi{\arabic{enumi}.}
\tightlist
\item
  В открытый контур входят, если специальный риск не доказан:

  \begin{enumerate}
  \def\labelenumii{\arabic{enumii}.}
  \tightlist
  \item
    законы, судебные решения, бюджеты, государственные контракты и
    данные о деятельности публичных органов;
  \item
    результаты исследований, преимущественно оплаченных из общественных
    средств;
  \item
    экологические, эпидемиологические, климатические и аварийные данные;
  \item
    стандарты безопасности, совместимости, доступности и ремонта;
  \item
    образовательные материалы и знания, необходимые для базовой
    медицины, воды, продовольствия, энергии и жизнеобеспечения;
  \item
    исторические и культурные материалы, срок ограниченного доступа к
    которым истёк;
  \item
    выводы, методика и сведения, необходимые обществу для оценки
    исполнения Соглашения, тогда как операционные подробности и
    защищаемые персональные данные передаются уполномоченному органу и
    публикуются только в безопасной форме.
  \end{enumerate}
\item
  Открытое знание публикуется в машиночитаемых и человекочитаемых
  форматах с указанием происхождения, версии, условий использования и
  известных ограничений.
\item
  Открытые условия предоставляют недискриминационное право законно
  копировать, переводить, адаптировать, зеркалировать, проверять,
  модифицировать, ремонтировать, обеспечивать совместимость, выполнять
  обратную разработку и применять знание в некоммерческом и коммерческом
  производстве с указанием происхождения и справедливым вознаграждением,
  когда оно установлено законом или лицензией.
\item
  Доступ к цифровой копии для изучения и проверки по общему правилу
  бесплатен. Разумная плата за физический носитель, производство, сервис
  или обеспеченное законом вознаграждение не должна создавать
  дискриминационный барьер.
\item
  Для новых соглашений об общественном финансировании открытая лицензия
  включается в условия заранее. Материал открывается при публикации либо
  не позднее двенадцати месяцев после принятия итогового отчёта
  финансирующим органом --- в зависимости от того, что наступит раньше,
  если суд не разрешил узкое ограничение по статье 21. Ранее возникшие
  права переводятся в открытый режим по статье 26 с соблюдением закона,
  договоров и справедливой компенсации.
\end{enumerate}

\subsubsection{Статья 21. Операционная информация высокого
риска}\label{ux441ux442ux430ux442ux44cux44f-21-ux43eux43fux435ux440ux430ux446ux438ux43eux43dux43dux430ux44f-ux438ux43dux444ux43eux440ux43cux430ux446ux438ux44f-ux432ux44bux441ux43eux43aux43eux433ux43e-ux440ux438ux441ux43aux430}

\begin{enumerate}
\def\labelenumi{\arabic{enumi}.}
\tightlist
\item
  Ограничение этого контура допустимо только при одновременном
  выполнении четырёх условий:

  \begin{enumerate}
  \def\labelenumii{\arabic{enumii}.}
  \tightlist
  \item
    сведения достаточно конкретны и непосредственно применимы;
  \item
    вред может быть тяжким, массовым или трансграничным;
  \item
    причинная связь между распространением и ростом риска обоснована
    проверяемыми фактами;
  \item
    редактирование, задержка, безопасная среда исследования или иной
    менее ограничительный режим недостаточны.
  \end{enumerate}
\item
  К возможным примерам относятся действующие ключи критической
  инфраструктуры, точные данные текущего наведения, актуальные планы
  уязвимых объектов и непосредственно исполнимые параметры причинения
  массового биологического, химического или ядерного вреда.
\item
  Само упоминание криптографии, искусственного интеллекта, биологии,
  полупроводников, робототехники, энергетики, космоса, оружия или
  двойного назначения не доказывает высокий риск.
\item
  Ограничение имеет дату окончания не позднее двух лет. Каждое продление
  требует нового решения суда, плана деклассификации и актуальных
  достаточных доказательств продолжающейся необходимости; строгость
  проверки возрастает, а после второго продления дело рассматривает
  коллегия не менее чем из трёх независимых судей.
\item
  Доверенные исследователи получают контролируемый доступ, когда это
  необходимо для проверки безопасности и разработки защиты.
\item
  Комиссия публикует обезличенное обоснование, объём ограничения и дату
  пересмотра, не раскрывая саму опасную часть.
\end{enumerate}

\subsubsection{Статья 22. Персональные и общинные
сведения}\label{ux441ux442ux430ux442ux44cux44f-22-ux43fux435ux440ux441ux43eux43dux430ux43bux44cux43dux44bux435-ux438-ux43eux431ux449ux438ux43dux43dux44bux435-ux441ux432ux435ux434ux435ux43dux438ux44f}

\begin{enumerate}
\def\labelenumi{\arabic{enumi}.}
\tightlist
\item
  Обработка персональных сведений основывается на законности,
  определённой цели, минимизации, точности, безопасности и ограниченном
  сроке хранения.
\item
  Особой защите подлежат медицинские, генетические, биометрические и
  интимные сведения, данные о детях и потерпевших, местонахождение
  уязвимых лиц и информация, создающая непосредственную угрозу
  преследования.
\item
  Лицо вправе требовать исправления, ограничения дальнейшего
  распространения, деиндексирования по имени, прекращения
  алгоритмического усиления, псевдонимизации и компенсации.
\item
  Если независимый орган установил необходимость и преобладающую
  доказательственную, научную или историческую ценность материала,
  полное уничтожение заменяется минимально достаточным запечатанным
  хранением после отделения избыточных данных и с периодическим
  пересмотром.
\item
  Сакральное и секретное традиционное знание охраняется с участием
  соответствующей общины. Его сохранение не создаёт автоматического
  права внешнего доступа.
\item
  Действующие ключи, пароли и полномочия доступа не являются
  цивилизационной памятью и подлежат безопасному отзыву или уничтожению
  после компрометации.
\item
  Стороны применяют открытые аудируемые протоколы защиты, минимизацию
  данных, сквозное шифрование, ротацию ключей и полномочий и
  воспроизводимые реализации там, где это технически уместно. Секретными
  являются ключи и конкретные полномочия доступа, а не необходимый для
  аудита алгоритм безопасности.
\end{enumerate}

\subsubsection{Статья 23. Факт публикации и эффективные средства
защиты}\label{ux441ux442ux430ux442ux44cux44f-23-ux444ux430ux43aux442-ux43fux443ux431ux43bux438ux43aux430ux446ux438ux438-ux438-ux44dux444ux444ux435ux43aux442ux438ux432ux43dux44bux435-ux441ux440ux435ux434ux441ux442ux432ux430-ux437ux430ux449ux438ux442ux44b}

\begin{enumerate}
\def\labelenumi{\arabic{enumi}.}
\tightlist
\item
  Средство защиты не основывается на юридической фикции гарантированного
  исчезновения всех копий уже широко опубликованного материала. Факт
  публикации устанавливается как факт и не означает законности материала
  или его дальнейшего распространения.
\item
  Невозможность гарантировать исчезновение всех копий не лишает
  пострадавшего эффективного средства защиты и не освобождает
  контролирующее распространение лицо от выполнимых обязанностей.
\item
  Средства защиты могут включать:

  \begin{enumerate}
  \def\labelenumii{\arabic{enumii}.}
  \tightlist
  \item
    прекращение контролируемого дальнейшего распространения;
  \item
    удаление из активных рекомендаций и поисковой выдачи по
    персональному идентификатору;
  \item
    исправление, контекстуализацию и обязательную связь с опровержением;
  \item
    редактирование персональных или опасных фрагментов;
  \item
    запечатанное хранение доказательственной копии;
  \item
    компенсацию и реабилитацию, а также передачу материалов
    компетентному органу для решения об ответственности за
    самостоятельное незаконное действие.
  \end{enumerate}
\item
  Средство должно быть технически выполнимым, адресным и соразмерным.
  Нельзя предписывать заведомо невозможное удаление «со всей планеты»
  или преследовать добросовестный независимый архив только за законное
  сохранение либо зеркалирование общественно значимого материала.
\item
  Ответственность определяется не только наличием информации, но и ролью
  лица, намерением, знанием риска, способом получения, способом усиления
  и фактически причинённым вредом.
\item
  Защита архива или посредника не распространяется на сознательное
  содействие тяжкому незаконному вреду, уклонение от адресной законной
  меры либо целенаправленное усиление материала при очевидном
  непосредственном риске.
\end{enumerate}

\subsubsection{Статья 24. Сохранение и запрет
уничтожения}\label{ux441ux442ux430ux442ux44cux44f-24-ux441ux43eux445ux440ux430ux43dux435ux43dux438ux435-ux438-ux437ux430ux43fux440ux435ux442-ux443ux43dux438ux447ux442ux43eux436ux435ux43dux438ux44f}

\begin{enumerate}
\def\labelenumi{\arabic{enumi}.}
\tightlist
\item
  Публичный орган, назначенный или профессиональный хранитель и лицо,
  которому вручено конкретное предписание о сохранении, не уничтожают
  заведомо уникальный объект существенной ценности либо объект,
  уникальность и ценность которого они разумно должны были знать,
  включая:

  \begin{enumerate}
  \def\labelenumii{\arabic{enumii}.}
  \tightlist
  \item
    последний известный экземпляр научно, культурно или исторически
    значимого материала;
  \item
    доказательства агрессии, основных международных преступлений и
    тяжких нарушений прав;
  \item
    официальные архивы, необходимые для подотчётности и прав
    потерпевших;
  \item
    данные безопасности, отсутствие которых препятствует предупреждению
    катастрофы;
  \item
    метаданные происхождения и историю изменений таких объектов.
  \end{enumerate}
\item
  Судебная палата или компетентный национальный суд может вынести
  срочное предписание о сохранении с немедленным пересмотром по
  ходатайству затронутого лица.
\item
  Обязанность не распространяется на каждого владельца каждой копии и не
  запрещает обычное управление дубликатами, повреждёнными данными,
  вредоносным кодом, скомпрометированными ключами и избыточными
  персональными сведениями.
\item
  Если открытый доступ создаёт несоразмерный вред, хотя сохранение
  необходимо, применяется риск-ориентированное запечатанное
  резервирование у организационно независимых хранителей в числе,
  достаточном для устойчивости и безопасности.
\item
  Если полный опасный объект невозможно сохранить безопасно, суд после
  независимой технической оценки может разрешить его проверяемое
  уничтожение, обязав сохранить безопасные репрезентативные данные, хэш,
  происхождение и доказательства принятого решения.
\item
  Сохранение не создаёт права доступа, а ограничение доступа не создаёт
  права уничтожить доказательство или уникальное знание.
\end{enumerate}

\subsubsection{Статья 25. Распределённые архивы, происхождение и
исправления}\label{ux441ux442ux430ux442ux44cux44f-25-ux440ux430ux441ux43fux440ux435ux434ux435ux43bux451ux43dux43dux44bux435-ux430ux440ux445ux438ux432ux44b-ux43fux440ux43eux438ux441ux445ux43eux436ux434ux435ux43dux438ux435-ux438-ux438ux441ux43fux440ux430ux432ux43bux435ux43dux438ux44f}

\begin{enumerate}
\def\labelenumi{\arabic{enumi}.}
\tightlist
\item
  Федерация цивилизационных архивов является сетью независимых хранилищ,
  а не единым мировым сервером или органом истины.
\item
  Архивы, где это безопасно и уместно, используют открытые форматы,
  контентную адресацию, криптографические подписи, проверяемую историю
  версий, географическое и организационное резервирование и
  офлайн-копии.
\item
  Хэш и подпись подтверждают целостность и заявленное происхождение, но
  не истинность содержания. Архив обязан ясно различать эти свойства.
\item
  Хранение, индексирование, рекомендация, научная оценка и политическое
  одобрение являются разными функциями и не приравниваются друг к другу.
\item
  Исправления, отзывы научных выводов, опровержения и существенные споры
  связываются с исходным объектом так, чтобы пользователь мог увидеть
  контекст, не уничтожая историческую запись.
\item
  Допускаются конкурирующие научные, культурные, образовательные и
  локальные индексы. Ни один индекс не получает исключительного права
  определять допустимое знание для всех.
\item
  Минимальные каталоги происхождения, открытых лицензий и мест хранения
  публикуются в совместимом формате.
\end{enumerate}

\subsubsection{Статья 26. Открытая наука, технологии и
вознаграждение}\label{ux441ux442ux430ux442ux44cux44f-26-ux43eux442ux43aux440ux44bux442ux430ux44f-ux43dux430ux443ux43aux430-ux442ux435ux445ux43dux43eux43bux43eux433ux438ux438-ux438-ux432ux43eux437ux43dux430ux433ux440ux430ux436ux434ux435ux43dux438ux435}

\begin{enumerate}
\def\labelenumi{\arabic{enumi}.}
\tightlist
\item
  Стороны переходят от технологической монополии к моделям, совмещающим
  открытый доступ, признание вклада и справедливое вознаграждение.
\item
  К новым общественно финансируемым исследованиям и стандартам
  критических гражданских систем применяется единый режим статьи 20. Для
  ранее возникших прав Стороны используют переходные лицензии, выкуп,
  компенсацию или согласованное открытие.
\item
  В чрезвычайной ситуации применяются обязательное лицензирование,
  патентные пулы, совместные закупки и передача производственных
  компетенций для лекарств, диагностики, воды, энергии, продовольствия и
  иных средств жизнеобеспечения.
\item
  Вознаграждение может обеспечиваться премиями, грантами, закупками,
  долгосрочными контрактами, открытыми консорциумами, оплатой
  производства и сервиса, репутационным признанием и переходными
  фондами.
\item
  Открытая лицензия не отменяет указание авторства, целостность научной
  атрибуции, трудовые права и согласованное справедливое распределение
  выгод.
\item
  В течение одного года принимается протокол об открытых гражданских
  технологиях для полупроводников, вычислений, криптографии, энергетики,
  медицины, материалов, робототехники, производства и космической
  инфраструктуры. Он устанавливает совместимость, права ремонта и
  обратной разработки, безопасную передачу производственных компетенций
  и защиту информации высокого риска.
\item
  В течение пяти лет Стороны устраняют необоснованные ограничения
  ремонта, совместимости, обратной разработки и воспроизводимости
  критических гражданских технологий.
\item
  Страны и сообщества с меньшими ресурсами получают финансирование
  инфраструктуры, переводов, образования и локального производства,
  чтобы формальная открытость стала фактической доступностью.
\end{enumerate}

\subsubsection{Статья 27. Общественно значимые
раскрытия}\label{ux441ux442ux430ux442ux44cux44f-27-ux43eux431ux449ux435ux441ux442ux432ux435ux43dux43dux43e-ux437ux43dux430ux447ux438ux43cux44bux435-ux440ux430ux441ux43aux440ux44bux442ux438ux44f}

\begin{enumerate}
\def\labelenumi{\arabic{enumi}.}
\tightlist
\item
  Лицо не подвергается ответственности или возмездию за общественно
  значимое раскрытие, если имело разумные основания считать сведения
  достоверными и общественно значимыми и приняло доступные разумные меры
  уменьшения несвязанного вреда.
\item
  Безопасные независимые каналы принимают материалы, защищают личность
  источника, проверяют подлинность и отделяют публичную часть от
  сведений высокого риска и персональных данных.
\item
  Предварительное обращение во внутренний канал не требуется, если он
  скомпрометирован, существует риск уничтожения доказательств, возмездия
  или непосредственного вреда.
\item
  Защита не распространяется на сознательную фабрикацию, вымогательство,
  насилие или намеренное раскрытие несвязанных сведений, когда тяжкий
  вред можно было разумно избежать.
\item
  Неправомерный способ доступа расследуется отдельно от общественной
  ценности раскрытого содержания; одно не предрешает другое.
\item
  Судебная палата может срочно запретить возмездие, сохранить
  доказательства и назначить безопасного хранителя.
\item
  Возмездием считается неблагоприятное действие, вызванное раскрытием
  или подготовкой к нему. Если такое действие последовало после
  защищённого раскрытия, обязанность доказать самостоятельное законное
  основание несёт применившее его лицо.
\end{enumerate}

\subsubsection{Статья 28. Посредники, фильтры и рекомендательные
системы}\label{ux441ux442ux430ux442ux44cux44f-28-ux43fux43eux441ux440ux435ux434ux43dux438ux43aux438-ux444ux438ux43bux44cux442ux440ux44b-ux438-ux440ux435ux43aux43eux43cux435ux43dux434ux430ux442ux435ux43bux44cux43dux44bux435-ux441ux438ux441ux442ux435ux43cux44b}

\begin{enumerate}
\def\labelenumi{\arabic{enumi}.}
\tightlist
\item
  Оператор, который добросовестно хранит, передаёт, кэширует или
  индексирует материал без целенаправленного содействия незаконному
  вреду, пользуется условной безопасной гаванью.
\item
  Безопасная гавань предполагает прозрачные правила, уведомление
  затронутых лиц, кроме случая конкретного риска для потерпевшего или
  доказательства, встречное уведомление, мотивированное решение, срочное
  обжалование и исполнение адресных законных мер, отвечающих статье 4 и
  независимому судебному контролю.
\item
  Оператор не обязан осуществлять общее наблюдение за всеми сообщениями
  или устанавливать универсальный фильтр.
\item
  Рекомендательная система несёт повышенную обязанность оценивать риск
  собственного алгоритмического усиления, особенно для прямого
  подстрекательства к насилию, преследования уязвимых лиц и материалов
  непосредственного массового вреда.
\item
  Каждый участник вправе использовать собственные фильтры и каталоги.
  Фильтр действует для добровольно выбравших его лиц и не становится
  глобальным правилом.
\item
  Публичный орган не может тайно требовать от посредника действий,
  которые он не вправе предписать открытым и оспоримым решением.
\item
  Коллизионные правила отдельного протокола определяют применимое право
  и разрешение противоречащих трансграничных предписаний так, чтобы одна
  юрисдикция не устанавливала безусловный глобальный режим.
\end{enumerate}

\subsection{Часть V. Институты
Соглашения}\label{ux447ux430ux441ux442ux44c-v-ux438ux43dux441ux442ux438ux442ux443ux442ux44b-ux441ux43eux433ux43bux430ux448ux435ux43dux438ux44f}

\subsubsection{Статья 29. Организация и Ассамблея
Сторон}\label{ux441ux442ux430ux442ux44cux44f-29-ux43eux440ux433ux430ux43dux438ux437ux430ux446ux438ux44f-ux438-ux430ux441ux441ux430ux43cux431ux43bux435ux44f-ux441ux442ux43eux440ux43eux43d}

\begin{enumerate}
\def\labelenumi{\arabic{enumi}.}
\tightlist
\item
  Настоящим Соглашением учреждается Организация планетарного мира с
  международной правосубъектностью и только той переданной компетенцией,
  которая прямо установлена Соглашением и ратифицированными протоколами.
\item
  Ассамблея определяет политику исполнения, принимает бюджет, избирает
  независимых должностных лиц, утверждает технические стандарты и
  проводит обзорные конференции, не присваивая полномочия национальных и
  иных международных органов.
\item
  Каждая государственная Сторона имеет один голос. Региональная
  организация голосует только вместо своих государств-членов в пределах
  заявленной компетенции и не увеличивает общее число их голосов.
\item
  Региональная организация не учитывается отдельно при определении
  кворума или порога в области, где голосуют её государства-члены.
\item
  Кворум составляет две трети всех учитываемых Сторон. Существенные
  решения принимаются не менее чем двумя третями всех учитываемых Сторон
  после разумной попытки достичь консенсуса; бюджет, поправка и адресная
  мера могут требовать более высокого порога по Соглашению.
\item
  Постоянного права вето нет.
\item
  Заседания, проекты решений, голоса, бюджеты и конфликты интересов
  открыты, кроме узко определённых приложений, защита которых необходима
  по статьям 21 или 22.
\end{enumerate}

\subsubsection{Статья 30. Совет предотвращения
конфликтов}\label{ux441ux442ux430ux442ux44cux44f-30-ux441ux43eux432ux435ux442-ux43fux440ux435ux434ux43eux442ux432ux440ux430ux449ux435ux43dux438ux44f-ux43aux43eux43dux444ux43bux438ux43aux442ux43eux432}

\begin{enumerate}
\def\labelenumi{\arabic{enumi}.}
\tightlist
\item
  Совет ведёт раннее предупреждение, поддерживает экстренную связь,
  предлагает посредничество и координирует краткосрочные меры
  предотвращения эскалации.
\item
  Его состав избирается на ограниченный срок с географическим, гендерным
  и правовым разнообразием; немедленное переизбрание ограничивается.
\item
  Совет может на срок не более семидесяти двух часов принять только
  невоенную, несиловую и временную координационную меру предотвращения
  эскалации. Для Сторон она рекомендательна; внутреннее поручение органу
  Организации обязательно только в пределах его существующей
  компетенции. Мера, затрагивающая право лица или Стороны,
  незамедлительно передаётся Судебной палате для проверки законности,
  необходимости и соразмерности и прекращается без её подтверждения.
\item
  Совет не вправе применять силу, задерживать лиц, проводить обыск или
  изъятие, назначать уголовное наказание, окончательно ограничивать
  информацию, вводить бессрочную секретность или распоряжаться
  уничтожением цивилизационной памяти.
\item
  Особое мнение члена Совета публикуется вместе с решением, если это не
  раскрывает защищаемые сведения.
\end{enumerate}

\subsubsection{Статья 31. Инспекция мира и
разоружения}\label{ux441ux442ux430ux442ux44cux44f-31-ux438ux43dux441ux43fux435ux43aux446ux438ux44f-ux43cux438ux440ux430-ux438-ux440ux430ux437ux43eux440ux443ux436ux435ux43dux438ux44f}

\begin{enumerate}
\def\labelenumi{\arabic{enumi}.}
\tightlist
\item
  Независимая инспекция проверяет декларации вооружений, меры снижения
  риска, защищённость опасных материалов и соблюдение запретов на атаки
  против охраняемых объектов.
\item
  Инспекторы действуют по профессиональным критериям, раскрывают
  конфликт интересов, не получают указаний от проверяемой Стороны и
  несут обязанность конфиденциальности.
\item
  Проверка ограничивается предметом мандата и не используется для
  промышленного, политического или военного шпионажа.
\item
  Публичный отчёт содержит методику, установленные факты, степень
  уверенности и рекомендации. Операционные детали, не необходимые
  обществу для оценки вывода, помещаются в защищённое приложение.
\item
  Ложное сообщение инспектора, сокрытие существенного конфликта
  интересов или передача защищённых данных влечёт независимое
  расследование и ответственность.
\item
  Обязательный проверочный протокол, принимаемый в течение ста
  восьмидесяти дней, регулирует уведомление и внезапность, управляемый
  доступ, отбор и хранение проб, конфиденциальность, привилегии
  инспекторов, мотивированный отказ, срочное разрешение споров и защиту
  от раскрытия несвязанных сведений.
\end{enumerate}

\subsubsection{Статья 32. Комиссия по знанию, правам и технологическому
риску}\label{ux441ux442ux430ux442ux44cux44f-32-ux43aux43eux43cux438ux441ux441ux438ux44f-ux43fux43e-ux437ux43dux430ux43dux438ux44e-ux43fux440ux430ux432ux430ux43c-ux438-ux442ux435ux445ux43dux43eux43bux43eux433ux438ux447ux435ux441ux43aux43eux43cux443-ux440ux438ux441ux43aux443}

\begin{enumerate}
\def\labelenumi{\arabic{enumi}.}
\tightlist
\item
  Комиссия разрабатывает открытые критерии трёх информационных контуров,
  независимо оценивает ограничения, рекомендует их подтверждение,
  изменение или отмену, координирует безопасное раскрытие и проверяет
  доступность общественно значимого знания.
\item
  В неё входят специалисты по правам человека, безопасности, науке,
  архивному делу, защите данных, традиционному знанию и технологиям из
  разных регионов.
\item
  Комиссия не является органом окончательной истины или глобальным
  цензором и не вводит окончательное обязательное ограничение. Её
  рекомендация мотивируется и публикуется в безопасной форме; решение о
  продлении принимает Судебная палата или компетентный национальный суд.
\item
  Научная неопределённость, конкурирующие оценки и особые мнения
  раскрываются явно.
\item
  Ежегодный отчёт показывает объём открытых материалов, действующие
  ограничения, сроки их пересмотра, исправления, жалобы и случаи
  уничтожения или утраты уникальной памяти.
\end{enumerate}

\subsubsection{Статья 33. Судебная палата и право
жалобы}\label{ux441ux442ux430ux442ux44cux44f-33-ux441ux443ux434ux435ux431ux43dux430ux44f-ux43fux430ux43bux430ux442ux430-ux438-ux43fux440ux430ux432ux43e-ux436ux430ux43bux43eux431ux44b}

\begin{enumerate}
\def\labelenumi{\arabic{enumi}.}
\tightlist
\item
  Судебная палата является неуголовным судебным органом, разрешает споры
  о толковании и применении Соглашения, рассматривает жалобы на решения
  его институтов и назначает эффективные средства защиты.
\item
  Межгосударственный спор может передать любая затронутая Сторона после
  этапов статьи 38, кроме просьбы о временной мере. Лицо или сообщество
  вправе после использования доступных эффективных внутренних средств
  либо при их очевидной недоступности подать жалобу против Стороны о
  нарушении охраняемого статьями 4--28 права или против института
  Соглашения о незаконном решении, непосредственно затронувшем
  заявителя.
\item
  При угрозе жизни, уничтожения доказательств, необратимой военной
  эскалации или раскрытия защищаемых сведений Палата принимает временные
  меры.
\item
  Разбирательство по общему правилу открыто. Защищённые доказательства
  рассматриваются в минимально необходимой закрытой части с
  процессуальными гарантиями для другой стороны.
\item
  Решения первой инстанции мотивируются, публикуются, имеют обязательную
  силу для участников дела и могут быть обжалованы в апелляционную
  коллегию Палаты по вопросам права и существенной фактической ошибки.
  Состав, независимость, отвод, сроки и исполнение определяются
  институциональным статутом, принимаемым не позднее ста восьмидесяти
  дней.
\end{enumerate}

\subsubsection{Статья 34. Федерация архивов и Фонд
перехода}\label{ux441ux442ux430ux442ux44cux44f-34-ux444ux435ux434ux435ux440ux430ux446ux438ux44f-ux430ux440ux445ux438ux432ux43eux432-ux438-ux444ux43eux43dux434-ux43fux435ux440ux435ux445ux43eux434ux430}

\begin{enumerate}
\def\labelenumi{\arabic{enumi}.}
\tightlist
\item
  Федерация цивилизационных архивов координирует совместимость и
  резервирование, не владея исключительным каталогом или единственной
  копией.
\item
  Членство открыто для соответствующих требованиям публичных,
  академических, общинных и частных хранилищ.
\item
  Фонд мира, восстановления и открытого знания финансирует разоружение,
  помощь потерпевшим, экологическое восстановление, безопасное
  уничтожение опасных материалов, открытые гражданские технологии,
  переводы, образование и локальные архивы.
\item
  Базовые взносы определяются по объективной шкале платёжеспособности с
  учётом дохода, населения и военных расходов. Исторически причинённый
  вред рассматривается отдельно как вопрос ответственности и репараций
  по решению компетентного органа. Помощь распределяется по прозрачным
  критериям потребности и результата.
\item
  Аудит Фонда публикует получателей, цели, контрольные точки, конфликты
  интересов и возврат средств при злоупотреблении, не раскрывая данные
  потерпевших.
\end{enumerate}

\subsection{Часть VI. Проверка, соблюдение и разрешение
споров}\label{ux447ux430ux441ux442ux44c-vi-ux43fux440ux43eux432ux435ux440ux43aux430-ux441ux43eux431ux43bux44eux434ux435ux43dux438ux435-ux438-ux440ux430ux437ux440ux435ux448ux435ux43dux438ux435-ux441ux43fux43eux440ux43eux432}

\subsubsection{Статья 35. Первоначальные декларации и национальное
исполнение}\label{ux441ux442ux430ux442ux44cux44f-35-ux43fux435ux440ux432ux43eux43dux430ux447ux430ux43bux44cux43dux44bux435-ux434ux435ux43aux43bux430ux440ux430ux446ux438ux438-ux438-ux43dux430ux446ux438ux43eux43dux430ux43bux44cux43dux43eux435-ux438ux441ux43fux43eux43bux43dux435ux43dux438ux435}

\begin{enumerate}
\def\labelenumi{\arabic{enumi}.}
\tightlist
\item
  В течение ста восьмидесяти дней после вступления Соглашения в силу для
  Стороны она представляет первоначальные декларации о вооружениях,
  опасных запасах, критических объектах, действующих информационных
  ограничениях и мерах защиты цивилизационной памяти.
\item
  Декларация может иметь публичную и защищённую части; объём закрытия
  обосновывается.
\item
  Каждая Сторона назначает независимый национальный орган координации,
  судебный путь обжалования и контакт экстренной связи.
\item
  Национальные законы устанавливают ответственность за
  воспрепятствование инспекции, уничтожение доказательств, ложную
  декларацию и возмездие против защищённого информатора.
\item
  Отсутствие ресурсов влечёт план помощи и срок исправления, а не
  автоматическое наказание.
\item
  Декларации и доступ к объектам осуществляются по проверочному
  протоколу статьи 31, включая управляемый доступ и защиту законных
  интересов, не связанных с предметом проверки.
\end{enumerate}

\subsubsection{Статья 36. Проверка и
доказательства}\label{ux441ux442ux430ux442ux44cux44f-36-ux43fux440ux43eux432ux435ux440ux43aux430-ux438-ux434ux43eux43aux430ux437ux430ux442ux435ux43bux44cux441ux442ux432ux430}

\begin{enumerate}
\def\labelenumi{\arabic{enumi}.}
\tightlist
\item
  Проверка использует сочетание регулярных и внезапных инспекций,
  дистанционных датчиков, спутниковых наблюдений, открытых источников,
  документов и показаний.
\item
  Каждый метод проходит оценку точности, уязвимости к подделке,
  дискриминационных последствий и защиты частной жизни.
\item
  Цепочка хранения доказательства, преобразования данных и
  использованные модели документируются так, чтобы независимый
  специалист мог проверить вывод.
\item
  Анонимный сигнал может начать проверку, но сам по себе не доказывает
  нарушение.
\item
  Проверяемая Сторона получает существенные факты и возможность ответить
  до окончательного вывода, кроме срочного риска уничтожения
  доказательств.
\item
  Наблюдение за частной перепиской населения не является допустимым
  общим способом проверки открытости знания или мирных обязательств.
\item
  Инспекция соблюдает проверочный протокол, а спор об отказе, объёме
  доступа, пробе или конфиденциальности разрешается Судебной палатой в
  срочном порядке.
\end{enumerate}

\subsubsection{Статья 37. Реакция на
несоблюдение}\label{ux441ux442ux430ux442ux44cux44f-37-ux440ux435ux430ux43aux446ux438ux44f-ux43dux430-ux43dux435ux441ux43eux431ux43bux44eux434ux435ux43dux438ux435}

\begin{enumerate}
\def\labelenumi{\arabic{enumi}.}
\tightlist
\item
  Реакция направлена на прекращение нарушения, уменьшение риска,
  восстановление прав и гарантии неповторения.
\item
  После уведомления, доступа к существенным материалам и возможности
  ответить Ассамблея большинством в три четверти всех учитываемых Сторон
  может применить следующую последовательность политических,
  экономических и организационных мер, а Судебная палата --- назначить
  судебные средства защиты в пределах своей компетенции:

  \begin{enumerate}
  \def\labelenumii{\arabic{enumii}.}
  \tightlist
  \item
    техническая помощь и согласованный план исправления;
  \item
    публичное предупреждение и усиленный мониторинг;
  \item
    временное приостановление добровольного финансирования, доступа к
    закупкам, техническим услугам протоколов или права занимать выборную
    должность в Организации;
  \item
    адресные экономические и организационные меры против определённых
    органов и лиц, ответственных за нарушение, только в пределах
    контроля Организации либо совместимого с международным правом
    национального исполнения;
  \item
    передача компетентному судебному или международному органу.
  \end{enumerate}
\item
  Срочный и умышленный тяжкий вред допускает переход непосредственно к
  необходимой ступени с последующим независимым пересмотром.
\item
  Меры не включают агрессивную силу, коллективное наказание, уничтожение
  гражданской инфраструктуры или сокрытие доказательств.
\item
  Не приостанавливаются права человека, гуманитарная помощь, доступ к
  правосудию, инспекция, сохранение доказательств и обязанности по
  безопасной ликвидации опасных материалов.
\item
  Адресат вправе обжаловать меру в Судебной палате; жалоба не
  приостанавливает срочную меру автоматически, но Палата может сделать
  это временно.
\item
  Выполнение плана исправления влечёт отмену допускающих отмену мер.
  Ответственность за уже причинённый вред сохраняется.
\end{enumerate}

\subsubsection{Статья 38. Порядок разрешения
споров}\label{ux441ux442ux430ux442ux44cux44f-38-ux43fux43eux440ux44fux434ux43eux43a-ux440ux430ux437ux440ux435ux448ux435ux43dux438ux44f-ux441ux43fux43eux440ux43eux432}

\begin{enumerate}
\def\labelenumi{\arabic{enumi}.}
\tightlist
\item
  Стороны сначала проводят консультации не более шестидесяти дней, если
  срочность не требует немедленной временной меры.
\item
  Затем спор передаётся посреднику или комиссии установления фактов на
  срок не более ста двадцати дней.
\item
  Неразрешённый спор передаётся в Судебную палату, если его участники
  совместно не согласовали арбитраж или иной обязательный форум.
  Производство не начинается и прекращается при наличии ранее начатого
  дела между теми же участниками, о том же предмете и на том же
  основании; окончательное решение исключает повторное рассмотрение.
\item
  Участники сохраняют доказательства, воздерживаются от ухудшения
  положения и выполняют временные меры.
\item
  Мирное соглашение по спору публикуется в части, необходимой для
  подотчётности, и не может нарушать неотъемлемые права третьих лиц.
\end{enumerate}

\subsubsection{Статья 39. Чрезвычайные
меры}\label{ux441ux442ux430ux442ux44cux44f-39-ux447ux440ux435ux437ux432ux44bux447ux430ux439ux43dux44bux435-ux43cux435ux440ux44b}

\begin{enumerate}
\def\labelenumi{\arabic{enumi}.}
\tightlist
\item
  При непосредственной угрозе массовой гибели, неконтролируемой
  эскалации или катастрофы Совет может принять координационную меру
  статьи 30. Судебная палата принимает обязательную временную меру по
  статье 33 только при непосредственном риске необратимого вреда
  охраняемому праву, раскрытия защищаемых сведений или уничтожения
  уникального доказательства.
\item
  Мера содержит фактическое основание, точный предмет, ответственного
  исполнителя, дату окончания и способ обжалования.
\item
  Мера Палаты действует не более тридцати дней, а каждое продление
  требует новых обстоятельств и не может увеличить совокупный срок свыше
  девяноста дней без полного разбирательства.
\item
  Чрезвычайная мера не разрешает пытку, рабство, насильственное
  исчезновение, дискриминацию, агрессию, безусловную цензуру или
  уничтожение доказательств. Она не создаёт самостоятельного основания
  для силы, задержания, обыска, изъятия, конфискации или уголовного
  наказания.
\item
  После прекращения угрозы публикуется отчёт о необходимости,
  результатах, побочном вреде и компенсации ошибочно затронутым лицам.
\end{enumerate}

\subsection{Часть VII. Переход и заключительные
положения}\label{ux447ux430ux441ux442ux44c-vii-ux43fux435ux440ux435ux445ux43eux434-ux438-ux437ux430ux43aux43bux44eux447ux438ux442ux435ux43bux44cux43dux44bux435-ux43fux43eux43bux43eux436ux435ux43dux438ux44f}

\subsubsection{Статья 40. Немедленные
меры}\label{ux441ux442ux430ux442ux44cux44f-40-ux43dux435ux43cux435ux434ux43bux435ux43dux43dux44bux435-ux43cux435ux440ux44b}

\begin{enumerate}
\def\labelenumi{\arabic{enumi}.}
\tightlist
\item
  Подписант может в пределах своего конституционного порядка временно
  применять всё Соглашение или перечисленные положения только
  посредством письменной декларации депозитарию. Декларация охватывает,
  в частности:

  \begin{enumerate}
  \def\labelenumii{\arabic{enumii}.}
  \tightlist
  \item
    отказ от нападений на гражданское население и усиленно защищённые
    объекты;
  \item
    гуманитарный доступ и экстренную связь;
  \item
    запрет уничтожения доказательств и уникальных архивов;
  \item
    уведомление об опасных инцидентах и ошибочных пусках;
  \item
    мораторий на новые системы смертельной автономии без значимого
    человеческого контроля;
  \item
    защиту лиц, имеющих разумные основания сообщать о непосредственной
    угрозе жизни;
  \item
    технологическое перемирие для критических гражданских средств
    жизнеобеспечения: временное обязательное лицензирование или
    патентный пул и неприменение запретительных исков, кроме адресной
    меры против доказанного высокого риска, нарушения частной жизни или
    опасного товара.
  \end{enumerate}
\item
  Временное применение не требует одномоментно публиковать действующие
  секреты, ключи или операционные сведения высокого риска.
\item
  Подписант может прекратить временное применение письменным
  уведомлением за тридцать дней. Прекращение не затрагивает
  ответственность за ранее совершённые действия, сохранение
  доказательств и приобретённые средства защиты.
\item
  Положение, требующее действия Инспекции, Судебной палаты или иного ещё
  не учреждённого института, временно применяется только после создания
  соответствующей переходной коллегии и прямого письменного принятия
  подписантом её компетенции.
\end{enumerate}

\subsubsection{Статья 41. Переходный
график}\label{ux441ux442ux430ux442ux44cux44f-41-ux43fux435ux440ux435ux445ux43eux434ux43dux44bux439-ux433ux440ux430ux444ux438ux43a}

\begin{enumerate}
\def\labelenumi{\arabic{enumi}.}
\tightlist
\item
  В первые сто восемьдесят дней Сторона создаёт национальный орган,
  подаёт декларации, инвентаризирует опасные материалы и критические
  архивы, проверяет экстренную связь и защищает каналы раскрытия.
  Ассамблея одновременно учреждает временные составы Судебной палаты и
  Инспекции, принимает институциональный статут и проверочный протокол;
  до этого их функции исполняют независимые переходные коллегии,
  избранные учредительной конференцией.
\item
  В течение двух лет она:

  \begin{enumerate}
  \def\labelenumii{\arabic{enumii}.}
  \tightlist
  \item
    приводит законы и практику в соответствие с Соглашением;
  \item
    открывает общественно финансируемые знания и стандарты критических
    гражданских систем;
  \item
    создаёт безопасные центры раскрытия уязвимостей;
  \item
    запускает архивные зеркала и компенсационные механизмы;
  \item
    начинает согласованные сокращения вооружений.
  \end{enumerate}
\item
  В течение одного года принимаются обязательные протоколы о вооружениях
  и открытых гражданских технологиях. В течение пяти лет завершаются
  первый обязательный цикл разоружения и проверки, пересмотр избыточной
  секретности и ограничений ремонта и формирование устойчивой Федерации
  архивов.
\item
  Переход сопровождается оценкой риска, безопасной миграцией, обучением,
  переводом, компенсацией и помощью Сторонам с меньшими возможностями.
\item
  Невыполнение контрольной точки требует публичного плана исправления с
  проверяемыми сроками.
\end{enumerate}

\subsubsection{Статья 42. Финансирование и развитие
возможностей}\label{ux441ux442ux430ux442ux44cux44f-42-ux444ux438ux43dux430ux43dux441ux438ux440ux43eux432ux430ux43dux438ux435-ux438-ux440ux430ux437ux432ux438ux442ux438ux435-ux432ux43eux437ux43cux43eux436ux43dux43eux441ux442ux435ux439}

\begin{enumerate}
\def\labelenumi{\arabic{enumi}.}
\tightlist
\item
  Общие институты имеют предсказуемый многолетний бюджет, защищённый от
  зависимости от одного государства или частного донора.
\item
  Обусловленная помощь не может требовать политической лояльности,
  отказа от законной научной независимости или передачи исключительного
  контроля над общим знанием.
\item
  Закупки осуществляются конкурентно и открыто, с раскрытием конечных
  владельцев, конфликтов интересов и критериев выбора.
\item
  Приоритет получают проекты, одновременно уменьшающие риск войны,
  технологическую зависимость и экологический вред.
\item
  Сообщества, непосредственно затронутые проектом, участвуют в его
  проектировании, мониторинге и оценке.
\end{enumerate}

\subsubsection{Статья 43. Поправки и периодический
обзор}\label{ux441ux442ux430ux442ux44cux44f-43-ux43fux43eux43fux440ux430ux432ux43aux438-ux438-ux43fux435ux440ux438ux43eux434ux438ux447ux435ux441ux43aux438ux439-ux43eux431ux437ux43eux440}

\begin{enumerate}
\def\labelenumi{\arabic{enumi}.}
\tightlist
\item
  Поправку может предложить Сторона, любой институт Соглашения или
  коалиция ассоциированных участников, поддержанная не менее чем пятью
  Сторонами.
\item
  Проект публикуется на рабочих языках, проходит оценку влияния на мир,
  права, безопасность, знание и окружающую среду и обсуждается не менее
  ста восьмидесяти дней.
\item
  Поправка принимается тремя четвертями Сторон и вступает в силу только
  для принявшей её Стороны после сдачи её ратификационной грамоты и
  достижения порога в две трети Сторон. В отношениях с не принявшей
  поправку Стороной продолжает действовать прежний текст.
\item
  Поправка к институциональным, бюджетным или юрисдикционным нормам
  вступает в силу только после принятия всеми Сторонами и действует для
  всех одновременно.
\item
  Чисто техническое приложение может обновляться двумя третями всех
  Сторон и вступать в силу через сто восемьдесят дней после уведомления,
  если Сторона не заявила отказ. Этот порядок не меняет права, запреты,
  юрисдикцию, финансовые обязательства или полномочия институтов; новая
  Сторона при присоединении указывает принимаемую действующую редакцию
  технических приложений.
\item
  Поправка, разрешающая агрессию, геноцид, рабство, пытку или иное
  нарушение императивной нормы международного права, недопустима и не
  имеет юридической силы независимо от числа согласившихся Сторон.
\item
  Иное изменение защищённого ядра, включая ослабление запрета
  уничтожения доказательств или эффективного судебного средства защиты,
  не может быть принято по процедуре настоящей статьи и требует нового
  договора, принятого всеми затрагиваемыми Сторонами.
\item
  Обзорная конференция проводится каждые пять лет с полноценным участием
  науки, гражданского общества, потерпевших, молодёжи и коренных
  народов.
\end{enumerate}

\subsubsection{Статья 44. Подписание, присоединение и вступление в
силу}\label{ux441ux442ux430ux442ux44cux44f-44-ux43fux43eux434ux43fux438ux441ux430ux43dux438ux435-ux43fux440ux438ux441ux43eux435ux434ux438ux43dux435ux43dux438ux435-ux438-ux432ux441ux442ux443ux43fux43bux435ux43dux438ux435-ux432-ux441ux438ux43bux443}

\begin{enumerate}
\def\labelenumi{\arabic{enumi}.}
\tightlist
\item
  Соглашение открыто для подписания всеми государствами и компетентными
  региональными интеграционными организациями.
\item
  Оно подлежит ратификации, принятию или утверждению в соответствии с
  внутренним правом каждого подписанта.
\item
  Региональная организация при сдаче грамоты представляет и обновляет
  декларацию компетенции и распределения ответственности. Её грамота не
  учитывается дополнительно к грамотам государств-членов при определении
  порога вступления в силу.
\item
  Соглашение вступает в силу через девяносто дней после сдачи сороковой
  государственной ратификационной грамоты при условии участия государств
  не менее чем из пяти региональных групп Организации Объединённых
  Наций.
\item
  Для последующего участника оно вступает в силу через девяносто дней
  после сдачи его грамоты.
\item
  Временное применение осуществляется только по статье 40.
\end{enumerate}

\subsubsection{Статья 45. Оговорки и
заявления}\label{ux441ux442ux430ux442ux44cux44f-45-ux43eux433ux43eux432ux43eux440ux43aux438-ux438-ux437ux430ux44fux432ux43bux435ux43dux438ux44f}

\begin{enumerate}
\def\labelenumi{\arabic{enumi}.}
\tightlist
\item
  Оговорки к Соглашению не допускаются.
\item
  Сторона может сделать толковательное заявление, если оно не исключает
  и не изменяет юридическое действие какого-либо обязательства.
\item
  Заявление, возражение против него и отзыв регистрируются и публикуются
  депозитарием. Судебная палата определяет юридическую природу спорного
  заявления по его содержанию, а не названию.
\end{enumerate}

\subsubsection{Статья 46. Выход и сохраняющиеся
обязательства}\label{ux441ux442ux430ux442ux44cux44f-46-ux432ux44bux445ux43eux434-ux438-ux441ux43eux445ux440ux430ux43dux44fux44eux449ux438ux435ux441ux44f-ux43eux431ux44fux437ux430ux442ux435ux43bux44cux441ux442ux432ux430}

\begin{enumerate}
\def\labelenumi{\arabic{enumi}.}
\tightlist
\item
  Сторона может выйти из Соглашения письменным уведомлением депозитария.
  Выход вступает в силу через двадцать четыре месяца.
\item
  Уведомление содержит мотивы, меры защиты затронутых лиц, порядок
  сохранения архивов и исполнение текущих планов безопасной ликвидации
  опасных материалов.
\item
  Выход не прекращает ответственность за действия, совершённые ранее,
  исполнение обязательных решений по текущим делам, сохранение
  доказательств, возмещение вреда и обязанности, действующие независимо
  от Соглашения.
\item
  Выход не даёт права изъять переданные законные копии из распределённой
  цивилизационной памяти или уничтожить совместные проверочные данные.
\item
  Депозитарий незамедлительно уведомляет всех Сторон и созывает
  консультацию о рисках выхода.
\end{enumerate}

\subsubsection{Статья 47. Депозитарий, регистрация и аутентичные
тексты}\label{ux441ux442ux430ux442ux44cux44f-47-ux434ux435ux43fux43eux437ux438ux442ux430ux440ux438ux439-ux440ux435ux433ux438ux441ux442ux440ux430ux446ux438ux44f-ux438-ux430ux443ux442ux435ux43dux442ux438ux447ux43dux44bux435-ux442ux435ux43aux441ux442ux44b}

\begin{enumerate}
\def\labelenumi{\arabic{enumi}.}
\tightlist
\item
  Депозитарием является Генеральный секретарь Организации Объединённых
  Наций либо иной депозитарий, единогласно назначенный учредительной
  конференцией.
\item
  Соглашение регистрируется и публикуется в соответствии с применимыми
  международными процедурами.
\item
  Арабский, английский, испанский, китайский, русский и французский
  тексты имеют одинаковую силу.
\item
  Депозитарий поддерживает открытый реестр подписаний, ратификаций,
  заявлений, поправок и уведомлений о выходе.
\item
  Соглашение толкуется добросовестно в соответствии с обычным значением
  терминов в их контексте и в свете предмета и цели. Последующее
  соглашение, последующая практика и применимые нормы международного
  права учитываются в установленном международным правом порядке;
  подготовительные материалы служат дополнительным средством толкования.
\end{enumerate}

\subsection{Приложение А. Минимальный переходный
календарь}\label{ux43fux440ux438ux43bux43eux436ux435ux43dux438ux435-ux430-ux43cux438ux43dux438ux43cux430ux43bux44cux43dux44bux439-ux43fux435ux440ux435ux445ux43eux434ux43dux44bux439-ux43aux430ux43bux435ux43dux434ux430ux440ux44c}

Сроки исчисляются со вступления Соглашения в силу для соответствующей
Стороны, а для временно применяемого положения --- со вступления в силу
её письменной декларации. Обязанность прекратить наступательные операции
возникает только для участника вооружённого конфликта по статье 6.

\begin{itemize}
\tightlist
\item
  \textbf{24 часа:} действуют экстренные линии связи; переданы
  уведомления о непосредственных опасностях.
\item
  \textbf{7 дней:} защищены гуманитарные маршруты, архивы и
  доказательства; назначены контактные лица.
\item
  \textbf{30 дней:} опубликован перечень временно применяемых мер и
  создан безопасный канал раскрытия.
\item
  \textbf{180 дней:} поданы первоначальные декларации, созданы
  национальные органы и планы исправления.
\item
  \textbf{1 год:} начаты инспекции, открытие общественно финансируемых
  знаний и пилотные архивные зеркала.
\item
  \textbf{2 года:} изменены основные законы, действуют механизмы
  компенсации и первые этапы разоружения.
\item
  \textbf{5 лет:} завершён первый цикл разоружения, проверки, архивной
  федерации и судебной защиты.
\end{itemize}

\subsection{Приложение Б. Проверка ограничения
доступа}\label{ux43fux440ux438ux43bux43eux436ux435ux43dux438ux435-ux431-ux43fux440ux43eux432ux435ux440ux43aux430-ux43eux433ux440ux430ux43dux438ux447ux435ux43dux438ux44f-ux434ux43eux441ux442ux443ux43fux430}

Орган, предлагающий ограничить материал как операционную информацию
высокого риска, письменно отвечает на все вопросы:

\begin{enumerate}
\def\labelenumi{\arabic{enumi}.}
\tightlist
\item
  Какая точная часть материала непосредственно применима?
\item
  Какой конкретный тяжкий вред и в какой временной перспективе разумно
  предсказуем?
\item
  Какие проверяемые факты связывают распространение с существенным
  ростом риска?
\item
  Почему редактирование, задержка, контролируемая исследовательская
  среда, уведомление уязвимой стороны или иной адресный способ
  недостаточны?
\item
  Кто пострадает от ограничения доступа и как этот вред будет уменьшен?
\item
  Каков срок, критерий досрочного снятия и независимый орган
  обжалования?
\item
  Какая безопасная копия сохранит доказательственную, научную или
  историческую ценность?
\end{enumerate}

Отсутствие содержательного ответа хотя бы на один обязательный вопрос
влечёт отказ в ограничении или его сужение.

\subsection{Приложение В. Минимальные данные
проверки}\label{ux43fux440ux438ux43bux43eux436ux435ux43dux438ux435-ux432-ux43cux438ux43dux438ux43cux430ux43bux44cux43dux44bux435-ux434ux430ux43dux43dux44bux435-ux43fux440ux43eux432ux435ux440ux43aux438}

Каждая проверочная запись содержит:

\begin{itemize}
\tightlist
\item
  предмет и правовое основание проверки;
\item
  дату, место и ответственных лиц;
\item
  источник, метод получения и цепочку хранения данных;
\item
  исходные измерения либо проверяемую ссылку на них;
\item
  все существенные преобразования, модели и допущения;
\item
  оценку неопределённости и известные альтернативные объяснения;
\item
  отделение установленного факта от аналитического вывода;
\item
  меры защиты персональных и операционных сведений;
\item
  ответ проверяемой стороны;
\item
  вывод, особые мнения, средство обжалования и срок следующего
  пересмотра.
\end{itemize}

\subsection{Опорные
материалы}\label{ux43eux43fux43eux440ux43dux44bux435-ux43cux430ux442ux435ux440ux438ux430ux43bux44b}

\begin{itemize}
\tightlist
\item
  \href{../источники/меморандум-неотзываемое-знание/memorandum_neotzyvaemoe_znanie.tex}{Меморандум
  о неотзываемом знании --- сохранённый TeX-источник}
\item
  \href{../источники/меморандум-неотзываемое-знание/отчёт-о-происхождении.md}{Отчёт
  о происхождении меморандума}
\item
  \href{https://www.un.org/en/about-us/un-charter/full-text}{Устав
  Организации Объединённых Наций}
\item
  \href{https://www.ohchr.org/Documents/ProfessionalInterest/cescr.pdf}{Международный
  пакт об экономических, социальных и культурных правах}
\item
  \href{https://www.unesco.org/en/legal-affairs/recommendation-open-science}{Рекомендация
  ЮНЕСКО по открытой науке}
\item
  \href{https://docs.un.org/A/RES/2222\%28XXI\%29}{Договор о космосе ---
  резолюция Генеральной Ассамблеи ООН 2222 (XXI)}
\item
  \href{https://www.un.org/en/genocide-prevention/1948-convention}{Конвенция
  о предупреждении преступления геноцида и наказании за него}
\item
  \href{../../запрос.md}{исходный запрос}
\end{itemize}

\end{document}
